\section{Confirmatory results}\label{sec:confirm}

Table~\ref{tab:verdicts} lists every pre-registered criterion with the value obtained and the
verdict. The reference values of the exploratory phase (at $w\ge1$) were reproduced to every digit
printed in the pre-registration before the confirmatory run.

\begin{table}[t]
\centering
\small
\caption{The five pre-registered hypotheses, each run once after the freeze. $\lamtwo$: modulus of
the second eigenvalue of the synapse-flow chain on the giant SCC; $\rho$: spectral radius of the
signed map; PR: participation ratio; $\Dob(P^r)$: Dobrushin coefficient of the $r$-step chain.}
\label{tab:verdicts}
\begin{tabular}{@{}l>{\raggedright\arraybackslash}p{7.5cm}>{\raggedright\arraybackslash}p{3.4cm}l@{}}
\toprule
 & Criterion (frozen) & Value & Met \\
\midrule
H1 (i) & $\lamtwo$ at $w\ge5$ lies in $[0.90,0.97]$ & $0.996524$ & no \\
H1 (ii) & removing ascending, descending and sensory-ascending neurons gives $1-\lamtwo<0.002$ & $0.000190$ & yes \\
H1 (iii) & removing the motor neurons moves $\lamtwo$ by less than $0.01$ & $0.000726$ & yes \\
\textbf{H1} & & & \textbf{fail} \\
\addlinespace
H2 (a) & with the floor $F$ (median input, 342 synapses), $\rho<0.85$ & $0.83484$ & yes \\
H2 (b) & at least 10 of the top 20 modes have PR $>20$ & 19 of 20 & yes \\
H2 (c) & descriptive: a top-20 mode with more than half its power on lobula-plate or ellipsoid-body neurons & mode 1, 99.3\,\% on the ellipsoid body & reported \\
\textbf{H2} & & & \textbf{pass} \\
\addlinespace
H3 (1) & certified lower bound on $\Dob(P^{32})$ at $w\ge5$ exceeds $0.20$ & $0.9229$ & yes \\
H3 (2) & first checkpoint with lower bound $<0.05$ lies in $[48,96]$ & none up to 128 ($0.6695$ at 128) & no \\
\textbf{H3} & & & \textbf{fail} \\
\addlinespace
H4 & certified lower bound on $\Dob(P'^{32})$ over the confidence class at $c=0.70$ exceeds $0.10$ & no certificate ($-0.52$) & no \\
H4 (clause) & largest level $c\in[0.51,0.95]$ at which the bound at $r=32$ exceeds $0.10$ & none (best: $0.071$ at $c=0.51$) & reported \\
\textbf{H4} & & & \textbf{fail} \\
\addlinespace
H5 & $\rho$ changes by less than $0.02$ under the Pospisil convention & $0.000015$ & yes \\
H5 & the carriers of the top six modes are unchanged & same six neurons & yes \\
\textbf{H5} & sensitivity result & & \textbf{pass} \\
\bottomrule
\end{tabular}
\end{table}

\paragraph{H1: the gap at five synapses is an order of magnitude smaller than predicted.}
At $w\ge1$ the chain has $\lamtwo=0.94503$, a spectral gap of $0.055$. At $w\ge5$ the prediction
was a gap between $0.03$ and $0.10$; the observed value is $\lamtwo=0.996524$, a gap of $0.0035$.
The two knock-out predictions held: without the ascending, descending and sensory-ascending neurons
the giant SCC (153{,}784 neurons) is still two-sided, held together by 18 residual neck neurons, and
has $1-\lamtwo=0.00019$; without the motor neurons $\lamtwo$ moves by $0.0007$.

\paragraph{H2 and H5: the photoreceptor modes are a normalisation effect.}
The leading eigenvalues of the signed map come in near-equal pairs of opposite sign
($-0.9151$, $+0.9137$, $-0.9122$, $+0.9098$), each mode carried by a pair of R7/R8 photoreceptors
(participation ratio 2.0 to 2.4). These photoreceptors receive few synapses (the six carriers have
2 to 43 input synapses from partners of nonzero sign, against 63 to 174 outputs), so postsynaptic
normalisation turns a reciprocal pair into a two-neuron loop of gain close to one. With the input floored at the
median, $F=342$, the spectral radius falls to $0.83484$ and 19 of the 20 leading modes extend over
more than 20 neurons. The leading floored mode is carried by GABAergic ellipsoid-body ring neurons
(ER3d\_b and ER3p\_a, participation ratio 96.6, 99.3\,\% of its power on neurons whose primary
neuropil is the ellipsoid body), and it is the ninth mode of the unfloored map with the same
eigenvalue: the floor removes the photoreceptor loops above it and leaves it untouched. Six of the
next eight modes are ring-neuron modes as well (ER4d, ER3p\_a, ER3m and three on ER3w\_b) and two
are on the leg sensory neurons LgAG1; the remaining modes of the top 20 are on Kenyon cells,
antennal-lobe projection neurons and sensory neurons of the cord. The result is the same with the floor at
700 synapses (18 of 20 modes extended). Under the Pospisil sign convention the unfloored spectral
radius changes by $1.5\times10^{-5}$ and the same six neurons carry the top six modes (H5).

\paragraph{H3: certified slow mixing.}
At $w\ge5$ the certified bounds are $0.923\le\Dob(P^{32})\le0.967$ and
$0.669\le\Dob(P^{128})\le0.816$ (Figure~\ref{fig:depth}): after 128 synaptic steps two starting
neurons can still be told apart with total-variation distance at least $0.669$. The first criterion held and the second failed,
because the bound never dropped below $0.05$. This is consistent with H1: at $\lamtwo=0.9965$ a
deviation decays by a factor ten only after several hundred steps.

\paragraph{H4: the worst-case certificate does not survive the confidence class.}
All 123{,}954{,}262 synapse-partner rows inside $G$ matched the connection weights exactly, and none
had a confidence below $0.5$. On average 14.4\,\% of a neuron's output synapses (12.3\,\% weighted by
the stationary distribution) have $\min(\text{conf}_\text{pre},\text{conf}_\text{post})<0.7$. With
per-row radii of this size, the certified lower bound of Lemma~\ref{lem:robust} is $0.821$ at $r=2$
and $0.396$ at $r=4$ and gives no certificate from $r=8$ on (Table~\ref{tab:h4}). Even at $c=0.51$,
where on average 0.5\,\% of a neuron's synapses lie below the level, the bound is $0.071$ at $r=32$, short of the
pre-registered $0.10$. As the pre-registration requires us to state: the depth result is not
certifiably robust to reconstruction uncertainty at the 0.7 level, nor at any level tested.

\begin{table}[t]
\centering
\small
\caption{H4: certified lower bounds on $\Dob(P'^r)$ over the confidence class at level $c$, and over
the uniform class $C(0.05)$, with 100 witness rows on $G$. A negative value means that no certificate
is obtained. The last row is the nominal bound for $P$ itself.}
\label{tab:h4}
\begin{tabular}{@{}lrrrrrrr@{}}
\toprule
$r$ & 1 & 2 & 4 & 8 & 16 & 32 & 48 \\
\midrule
$c=0.51$ & 1.000 & 0.997 & 0.977 & 0.860 & 0.555 & 0.071 & $-0.015$ \\
$c=0.55$ & 1.000 & 0.984 & 0.894 & 0.595 & $-0.028$ & $-0.071$ & $-0.071$ \\
$c=0.60$ & 1.000 & 0.941 & 0.759 & 0.264 & $-0.192$ & $-0.198$ & $-0.199$ \\
$c=0.70$ & 1.000 & 0.821 & 0.396 & $-0.470$ & $-0.511$ & $-0.520$ & $-0.522$ \\
$C(0.05)$ & 0.789 & 0.579 & 0.156 & $-0.441$ & $-0.481$ & $-0.502$ & $-0.508$ \\
\addlinespace
nominal $P$ & 1.000 & 1.000 & 0.998 & 0.936 & 0.707 & 0.312 & 0.127 \\
\bottomrule
\end{tabular}
\end{table}
